\documentclass[11pt]{article}
\usepackage[margin=1in]{geometry}
\usepackage{amsmath,amssymb,amsthm}
\usepackage[T1]{fontenc}
\usepackage{lmodern}
\usepackage[hidelinks]{hyperref}
\newtheorem{theorem}{Theorem}
\newtheorem{lemma}{Lemma}
\newcommand{\Pp}{\mathbb P}
\newcommand{\Ee}{\mathbb E}
\newcommand{\Var}{\operatorname{Var}}
\newcommand{\R}{\mathbb R}
\newcommand{\Z}{\mathbb Z}
\title{Sharp modulus range for triangle counts\\in dense random graphs}
\author{Alec Kriebel\\
\small ORCID: \href{https://orcid.org/0009-0001-9320-500X}{0009-0001-9320-500X}}
\date{Version 1.0 --- 17 October 2026}
\begin{document}
\maketitle
\begin{abstract}
Fix \(0<p<1\), and let \(T_n\) count each triangle once in \(G(n,p)\).
For every positive integer sequence \(q_n\), we give the explicit deduction
\[
 \max_{0\leq a<q_n}
 \left|q_n\Pp(T_n\equiv a\pmod {q_n})-1\right|\longrightarrow0
 \quad\Longleftrightarrow\quad q_n=o(n^2).
\]
Sufficiency follows from the quantitative local limit theorem of Sah and
Sawhney, concentration, and finite-window summation. Necessity follows from
an elementary edge-projection central limit theorem and a modular Fourier
coefficient. This covers arbitrary, possibly oscillating or composite,
modulus sequences. We also reconcile density and Gaussian unit-cell
conventions. The result is a credited corollary of earlier local limit
work, addressing the sharp range in Scott's 2009 report of joint work with
Tateno. No novelty, historical firstness, or new local limit theorem is claimed.
\end{abstract}

\section{Question, conventions, and result}

Scott's report of joint work with Tateno \cite[pp.~1279--1280]{ST}
asks how far equidistribution of the triangle count modulo \(q\) extends.
It announces the relative, uniform-in-residue conclusion for
\(q=o(n^2/\log n)\), without a primality restriction, already identifies
\(o(n^2)\) as the necessary scale, and announces a triangle local limit
theorem. Here we spell out the complete sharp-range deduction from a later
quantitative local limit theorem. This is a present-day deduction from
prior results, not an assertion that this corollary was explicitly stated
in an earlier source or that the question remained open until this note.
The accompanying bounded literature audit is not a certificate of novelty.

The triangle local limit theorem was established by Gilmer and Kopparty
\cite{GK}, and quantitative pointwise estimates were obtained by Berkowitz
\cite{B}. These are substantive predecessors of the connected-count theorem
used below, rather than new local limit results of this note.

Throughout, \(p\in(0,1)\) is fixed, the edges of \(G(n,p)\) are independent,
and \(T_n\) is the number of unordered three-vertex sets spanning a triangle.
Thus each triangle is counted once; a labeled-embedding count would be
\(6T_n\) and would have different modular restrictions. Write
\[
 \mu_n=\Ee T_n,\qquad \sigma_n^2=\Var T_n,\qquad
 \Delta_n(q)=\max_{0\leq a<q}
       \left|q\Pp(T_n\equiv a\pmod q)-1\right|.
\]
All asymptotic statements are as \(n\to\infty\), and implicit constants may
depend on fixed \(p\). Standard-deviation normalization starts at \(n\geq3\).

\begin{theorem}\label{thm:main}
For every sequence of positive integers \(q_n\),
\[
 \Delta_n(q_n)\longrightarrow0
       \quad\Longleftrightarrow\quad q_n=o(n^2).
\]
No monotonicity, primality, or lower growth condition on \(q_n\) is required.
\end{theorem}

In particular \(q_n=1\) is permitted. Sparse or varying \(p\), general
subgraph counts, and induced-copy analogues are outside this theorem.
The broader final conjecture in \cite{ST} concerns arbitrary fixed
subgraphs and is a separate question. Its unrestricted Gaussian local
limit part already has the ordinary \(2K_2\) counterexample of Sah and
Sawhney \cite[Theorem~1.3]{SS}; that known obstruction is not a new result
of this note.

\section{Moments and the imported local theorem}

Let \(I_e\) be the edge indicators, \(\xi_e=I_e-p\), and
\(N=\binom n2\). Expanding the three factors for each triangle gives
\begin{align}
 T_n-\mu_n&=L_n+R_n,\qquad
 L_n=p^2(n-2)\sum_e\xi_e, \label{eq:projection}\\
 R_n&=p\sum_{\{e,f\}\ {\rm adjacent}}\xi_e\xi_f
         +\sum_{\{e,f,g\}\ {\rm triangle}}\xi_e\xi_f\xi_g. \notag
\end{align}
Each unordered adjacent edge pair occurs once. Distinct edge sets index
orthogonal centered products, since an edge in their symmetric difference
has a centered factor occurring only once. There are \(3\binom n3\)
adjacent pairs. Consequently, with \(\alpha=p^5(1-p)>0\) and
\(D=N(n-2)^2\),
\begin{align}
 \mu_n&=\binom n3p^3, \notag\\
 \sigma_n^2
 &=\alpha D+3\binom n3p^4(1-p)^2+\binom n3p^3(1-p)^3 \label{eq:variance}\\
 &=\binom n3p^3(1-p^3)+12\binom n4p^5(1-p)
       \sim\frac{\alpha}{2}n^4. \notag
\end{align}
The coefficient \(12\binom n4\) counts ordered pairs of distinct triangles
sharing an edge. Thus
\begin{equation}\label{eq:scale}
 \sigma_n^2\geq\alpha D,\qquad
 \sigma_n\sim\sqrt{\alpha/2}\,n^2.
\end{equation}
The empty graph and a graph consisting of exactly one triangle and isolated
vertices have positive probability, so the integer span is one. We do not
assume every intermediate count is attainable.

Our substantive external premise is Sah--Sawhney's quantitative
local limit theorem \cite[Theorem~1.1]{SS}, in the inspected
arXiv:2006.11369v1. For ordinary connected subgraph counts it uses the
actual mean and standard deviation in \(G(n,p)\). Its ordinary-copy
definition and independent-edge transfer are in Sections~4.1 and~5.
Specializing to \(K_3\), and choosing
\(\lambda=\min(p,1-p)/2\) in its density hypothesis, gives, for every
fixed \(\varepsilon>0\),
\begin{equation}\label{eq:llt}
 \sup_{k\in\Z}\left|\Pp(T_n=k)-g_n(k)\right|
 \leq\frac{\delta_n}{\sigma_n},\qquad
 g_n(x)=\frac1{\sigma_n}\phi\!\left(\frac{x-\mu_n}{\sigma_n}\right),
 \quad \delta_n=C_{p,\varepsilon}n^{-1/2+\varepsilon},
\end{equation}
where \(\phi(x)=(2\pi)^{-1/2}e^{-x^2/2}\). The theorem counts each
triangle once, rather than its six embeddings. We import this established,
deep Fourier/decoupling theorem; we do not reprove or independently certify
its entire proof. Only its uniform pointwise bound is needed. Fix
\(0<\varepsilon<1/2\), so \(\delta_n\sqrt{\log n}\to0\).
An unquantified pointwise error cannot simply be summed over the full
count support to obtain our relative modular assertion.

\section{Sufficiency: finite-window summation}

Changing one edge changes \(T_n\) by at most \(n-2\). The independent-edge
exposure martingale and the bounded-range exponential-moment inequality
give
\begin{equation}\label{eq:concentration}
 \Pp(|T_n-\mu_n|\geq u)\leq2e^{-2u^2/D}.
\end{equation}
Indeed, each martingale increment has conditional range length at most
\(n-2\); its centered exponential moment is at most
\(e^{t^2(n-2)^2/8}\). Multiplying these bounds, applying exponential
Markov, and optimizing \(t=4u/D\) proves each one-sided bound and hence
\eqref{eq:concentration}.

Set
\[
 B_n=\sqrt{(16/\alpha)\log n}.
\]
By \eqref{eq:scale} and \eqref{eq:concentration}, the probability outside
\(|T_n-\mu_n|\leq B_n\sigma_n\) is at most \(2n^{-32}\).
For large \(n\), \(\sigma_n\geq1\) and \(B_n\geq4\).
Comparing each Gaussian grid point on a tail with its adjacent interval
toward the origin gives
\begin{align}
 \sum_{\substack{k\in\Z\\|k-\mu_n|>B_n\sigma_n}}g_n(k)
 &\leq2\int_{B_n-1/\sigma_n}^{\infty}\phi(x)\,dx \notag\\
 &\leq2e^{-(B_n-1)^2/2}
 \leq2e^{-B_n^2/4}
 =2n^{-4/\alpha}\leq2n^{-16}. \label{eq:gausstail}
\end{align}
Here \(\alpha\leq1/4\), and
\((B-1)^2/2\geq B^2/4\) for \(B\geq4\). These are full tails, so they
bound each residue's omitted mass as well.

Fix \(1\leq q\leq\sigma_n\) and a residue \(a\). There are at most
\(2B_n\sigma_n/q+1\) integers in \(a+q\Z\) inside the central interval.
Use \eqref{eq:llt} only on those finitely many integers, and use the two
full tails elsewhere. We obtain
\begin{equation}\label{eq:window}
 q\left|\Pp(T_n\equiv a\pmod q)
                  -\sum_{k\in a+q\Z}g_n(k)\right|
 \leq\delta_n(2B_n+q/\sigma_n)+2q(n^{-32}+n^{-16}).
\end{equation}
All factors of \(q\) are retained.

We also use the elementary shifted-mesh inequality
\begin{equation}\label{eq:mesh}
 \left|h\sum_{j\in\Z}f(b+jh)-\int_{\R}f(x)\,dx\right|
       \leq h\int_{\R}|f'(x)|\,dx
       \qquad(h>0,\ b\in\R),
\end{equation}
for absolutely continuous \(f\) with \(f,f'\in L^1(\R)\).
To prove it, compare \(h f(b+jh)\) with the integral over
\([b+jh,b+(j+1)h]\); the difference is bounded by \(h\) times the
integral of \(|f'|\) over that interval. Summation proves the bound.
Applying the same comparison to \(|f|\) justifies absolute convergence.
For \(\phi\), \(\int\phi=1\) and
\(\int|\phi'|=\sqrt{2/\pi}\). Taking \(h=q/\sigma_n\) and
\(b=(a-\mu_n)/\sigma_n\) in \eqref{eq:mesh} yields
\[
 \left|q\sum_{k\in a+q\Z}g_n(k)-1\right|
           \leq\sqrt{2/\pi}\,q/\sigma_n .
\]
Together with \eqref{eq:window}, this proves the explicit uniform bound
\begin{equation}\label{eq:discrepancy}
 \Delta_n(q)\leq
 \delta_n(2B_n+q/\sigma_n)
 +2q(n^{-32}+n^{-16})+\sqrt{2/\pi}\,q/\sigma_n,
 \qquad 1\leq q\leq\sigma_n .
\end{equation}
If \(q_n=o(\sigma_n)\), the bound applies eventually and every term tends
to zero, since \(\sigma_n=O(n^2)\).
This proves sufficiency even when \(q_n/\sigma_n\) tends to zero
arbitrarily slowly.

\section{Necessity and the original cell convention}

For sharpness, the local limit theorem is unnecessary. From
\eqref{eq:projection}--\eqref{eq:variance},
\(\Var L_n=\alpha D=\Theta_p(n^4)\) and \(\Var R_n=O_p(n^3)\).
The binomial central limit theorem for \(\sum_e I_e\), together with
\(R_n/\sigma_n\to0\) in \(L^2\), gives
\[
 Z_n=(T_n-\mu_n)/\sigma_n\ \Longrightarrow\ N(0,1),
 \qquad \Ee|Z_n|\leq1 .
\]
Suppose \(\Delta_n(q_n)\to0\) but \(q_n/\sigma_n\not\to0\).
On an infinite subsequence \(q_n/\sigma_n\geq c>0\); passing to a further
subsequence, \(r_n=\sigma_n/q_n\to r\in[0,1/c]\).
Eventually \(q_n\geq2\), so the sum of the \(q_n\)-th roots of unity is zero
and
\begin{equation}\label{eq:fourier}
 \left|\Ee e^{2\pi iT_n/q_n}\right|
 \leq\sum_{a=0}^{q_n-1}
       \left|\Pp(T_n\equiv a\pmod {q_n})-1/q_n\right|
 \leq\Delta_n(q_n)\longrightarrow0 .
\end{equation}
Removing the centering phase leaves
\(\left|\Ee e^{it_nZ_n}\right|\), where \(t_n=2\pi r_n\to2\pi r=t\).
The varying frequency is legitimate because
\[
 \left|\Ee e^{it_nZ_n}-\Ee e^{itZ_n}\right|
       \leq |t_n-t|\,\Ee|Z_n|\leq|t_n-t|\longrightarrow0 .
\]
The central limit theorem therefore makes the left side of
\eqref{eq:fourier} tend to \(e^{-2\pi^2r^2}>0\), including \(1\) when
\(r=0\), a contradiction. Thus \(q_n=o(\sigma_n)\), equivalently
\(q_n=o(n^2)\). This subsequence argument includes nonmonotone, critical,
and supercritical sequences, and completes the proof of
Theorem~\ref{thm:main}.

For comparison with the original local theorem in \cite{ST}, define its
centered Gaussian unit-cell mass
\[
 c_n(k)=
 \int_{(k-\mu_n-1/2)/\sigma_n}^{(k-\mu_n+1/2)/\sigma_n}\phi(x)\,dx .
\]
Taylor's theorem about \((k-\mu_n)/\sigma_n\), integrated symmetrically,
cancels the linear term and proves, uniformly in \(k\),
\[
 |c_n(k)-g_n(k)|\leq\frac{\|\phi''\|_\infty}{24\sigma_n^3}.
\]
Hence \eqref{eq:llt} also gives
\[
 \sup_{k\in\Z}|\Pp(T_n=k)-c_n(k)|
 =O_{p,\varepsilon}(n^{-5/2+\varepsilon})+O_p(n^{-6})
 =o(n^{-2})\qquad(0<\varepsilon<1/2).
\]
Thus the density convention of the imported theorem also supplies the
original cell convention; this is not a change of the triangle target.

\paragraph{Verification and disclosure.}
The supporting materials distinguish fresh exact finite checks, proof
reviews, and source inspection from historical records. Finite enumeration
checks the moment and counting identities; it does not establish an
asymptotic theorem. AI tools were used extensively in developing,
reconstructing, drafting, and verifying this work, including independent
internal AI adversarial reviews. This note is unrefereed and has received
no external human peer review or formal proof-assistant certification.
The imported local limit theorem remains a named external mathematical
premise. No novelty or historical-priority claim is made.

\begin{thebibliography}{9}
\bibitem{GK}
J.~Gilmer and S.~Kopparty,
\emph{A local central limit theorem for triangles in a random graph},
Random Structures Algorithms \textbf{48} (2016), no.~4, 732--750.
\href{https://doi.org/10.1002/rsa.20604}{doi:10.1002/rsa.20604}.
Author manuscript dated 6 December 2014:
\href{https://sites.math.rutgers.edu/~sk1233/triangles.pdf}{author PDF}.

\bibitem{B}
R.~Berkowitz,
\emph{A Quantitative Local Limit Theorem for Triangles in Random Graphs},
\href{https://arxiv.org/abs/1610.01281v3}{arXiv:1610.01281v3}
(submitted 2016; archive revision 2017).

\bibitem{ST}
A.~Scott (joint work with A.~Tateno),
\emph{On the Number of Triangles in a Random Graph},
in \emph{Combinatorics and Probability}, Oberwolfach Report 23/2009,
pp.~1279--1280.
\href{https://doi.org/10.4171/OWR/2009/23}{doi:10.4171/OWR/2009/23}.

\bibitem{SS}
A.~Sah and M.~Sawhney,
\emph{Local limit theorems for subgraph counts},
J.~Lond.\ Math.\ Soc.\ (2) \textbf{105} (2022), no.~2, 950--1011.
\href{https://doi.org/10.1112/jlms.12523}{doi:10.1112/jlms.12523}.
Preprint version used:
\href{https://arxiv.org/abs/2006.11369v1}{arXiv:2006.11369v1} (2020).
\end{thebibliography}
\end{document}
